% META: {"id": "TEXP-ARI-CO-042", "chapitre": "TEXP-ARITHMETIQUE", "type_objet": "corrige", "exercice_ref": "TEXP-ARI-EX-042", "capacites_codes": ["C6"], "sources_inspiration": [], "status": "approved"}
% BEGIN-VERIFY
% from sympy import *
% x = symbols('x')
% # f(x) = x^4 - 4x^2 + 1, trouver nb de points d'inflexion et leur nature
% f = x**4 - 4*x**2 + 1
% fpp = diff(f, x, 2)
% assert expand(fpp) == 12*x**2 - 8
% assert solve(fpp, x) == [-sqrt(6)/3, sqrt(6)/3]
% assert fpp.subs(x, 0) == -8  # concave
% assert fpp.subs(x, 2) == 40  # convexe
% f.subs(x, sqrt(6)/3)
% assert simplify(f.subs(x, sqrt(6)/3) - Rational(-11, 9)) == 0
% END-VERIFY
\begin{corrige}{TEXP-ARI-EX-042}

\textbf{Question 1.} $f''(x) = 12x^2 - 8 = 4(3x^2-2)$. $f'' = 0 \iff x = \pm\dfrac{\sqrt{6}}{3}$.

\commentaireMarge{$f''=4(3x^2-2)$ : deux racines, deux points d'inflexion.}

Deux points d'inflexion : $f''$ change de signe en chaque racine.

\textbf{Question 2.} $f'(x) = 4x^3 - 8x = 4x(x^2-2)$, nul en $0, \pm\sqrt{2}$.

$f$ est convexe sur $\left]-\infty, -\frac{\sqrt6}{3}\right] \cup \left[\frac{\sqrt6}{3}, +\infty\right[$ et concave entre.

\textbf{Question 3.} $f(0)=1>0$, $f(\sqrt2) = 4-8+1 = -3 < 0$, $\lim_{\pm\infty} f = +\infty$.

Par continuite, la courbe coupe l'axe $4$ fois (TVI sur chaque intervalle monotone).

\end{corrige}
